%=============================================
\section{System Architecture}
%=============================================

\subsection{High-Level Architecture}
The system is built upon a microservices architecture, heavily utilizing containerization to isolate and manage dependencies. It consists of nine interconnected services communicating over a dedicated Docker bridge network. Figure~\ref{fig:arch} illustrates the complete system architecture, showing the interaction between the frontend, backend, stream processing pipeline, and the TTS engine.

\begin{figure}[H]
\centering
\includegraphics[width=\textwidth]{system_architecture_1777743769425.png}
\caption{High-level system architecture of the Real-Time AI Cricket Commentator}
\label{fig:arch}
\end{figure}

\subsection{Data Flow Pipeline}
The system operates on a continuous loop of event generation, processing, and feedback. The end-to-end data flow is illustrated in Figure~\ref{fig:dataflow} and proceeds as follows:

\begin{enumerate}
\item \textbf{Event Generation}: The user interacts with the React frontend, triggering a cricket game event (e.g., hitting a six, taking a wicket). The browser constructs a JSON payload detailing the event.
\item \textbf{Ingestion}: The browser sends the JSON event to the Node.js Backend API via an open WebSocket connection. The backend acts as a Kafka Producer, immediately publishing the event to the \texttt{game-events} Kafka topic.
\item \textbf{Stream Processing}: Apache Flink, acting as a Kafka Consumer, ingests the \texttt{game-events} topic. The data stream is keyed by \texttt{player\_id} and duplicated across five parallel processors (Kill Streak, Aggression, Inactivity, Momentum, Win Predictor).
\item \textbf{Commentary Output}: The Flink processors analyze the event within the context of recent history. If a significant pattern is detected, they generate a commentary trigger. These triggers are unified into a single stream and published to the \texttt{commentary} Kafka topic.
\item \textbf{Broadcasting}: The Node.js Backend consumes the \texttt{commentary} topic. It applies emotion tagging and formatting, then broadcasts the text commentary back to all connected React clients via WebSocket.
\item \textbf{Voice Synthesis}: Simultaneously, the backend makes an asynchronous HTTP POST request to the Python TTS service, requesting voice synthesis for the commentary string. The TTS service returns a Base64-encoded WAV file, which the backend then forwards to the browser via WebSocket for audio playback.
\end{enumerate}

\begin{figure}[H]
\centering
\includegraphics[width=\textwidth]{data_flow_pipeline_1777743783686.png}
\caption{Complete data flow pipeline from event generation to commentary delivery}
\label{fig:dataflow}
\end{figure}

\subsection{Technology Stack}
The project utilizes a modern, robust technology stack suitable for production-grade big data applications.

\begin{table}[H]
\centering
\caption{Technology Stack Overview}
\label{tab:techstack}
\begin{tabular}{@{}lll@{}}
\toprule
\textbf{Layer} & \textbf{Technology} & \textbf{Version} \\
\midrule
Event Streaming & Apache Kafka (Confluent) & 7.5.0 \\
Stream Processing & Apache Flink (DataStream API) & 1.18.1 \\
Game Simulator & HTML5 Canvas / React & 18.x \\
Backend API & Node.js + Express + WebSocket (ws) & 20.x \\
Frontend UI & React + Vite & 18 / 5 \\
TTS Engine & Python + Coqui XTTS-v2 & 3.11 \\
Containerization & Docker + Docker Compose & Latest \\
Reverse Proxy & Nginx & Alpine \\
Coordination & Apache Zookeeper & 7.5.0 \\
Build Tool & Apache Maven & 3.9 \\
\bottomrule
\end{tabular}
\end{table}

\subsection{Communication Protocols}
The system employs three distinct communication protocols optimized for their specific use cases:

\begin{figure}[H]
\centering
\includegraphics[width=\textwidth]{websocket_realtime_1777743907214.png}
\caption{WebSocket bidirectional communication between frontend and backend}
\label{fig:websocket}
\end{figure}

\begin{itemize}
\item \textbf{Apache Kafka (TCP)}: Used for asynchronous, high-throughput, persistent event streaming between the backend Node.js server and the Flink cluster. It decouples the web tier from the analytics tier.
\item \textbf{WebSocket (TCP)}: Used for real-time, low-latency, bidirectional communication between the user's browser and the backend server. It allows the server to push updates (commentary, stats, audio) to the client without the overhead of HTTP polling.
\item \textbf{REST API (HTTP)}: Used for synchronous communication between the backend server and the Python TTS service, and for general health checks.
\end{itemize}
